\section{Independent counts and uniform word order}
\label{sec:poisson_uniform_order}

Retain the finite alphabet $I$, the word law $\mu_t$ for $t\ge0$, and
$C_n=\mathcal W_n=\{w:C(w)=n\}$ for $n\in\N^I$.
The following is a characterization at a single fixed time, supplying the
count-and-order description relevant to \cite{OpenAI2026AreaLaw},
\texttt{09-amplification.tex}, lines 49--54 and 237--253.

\begin{theorem}[Conditional word order]
\label{thm:poisson_conditional_order}
\lean{PoissonWord.cond_countEvent, PoissonWord.measure_countEvent_pos,
PoissonWord.cond_countEvent_of_pos}
\leanok
\uses{def:poisson_word_measure, def:poisson_word_count_vector}
If $\mu_t(C_n)\ne0$, then the actual conditional word law is
\begin{align}
 \mu_t(\,\cdot\mid C_n)&=\operatorname{Unif}(C_n).
 \label{eq:poisson_conditional_order}
\end{align}
For $t>0$, every count vector $n$ satisfies $\mu_t(C_n)>0$, so this
conditional identity holds for every $n\in\N^I$.
\end{theorem}
\begin{proof}
\leanok
\uses{lem:poisson_word_count_fiber, lem:poisson_count_event_mass,
lem:poisson_word_normalization, thm:finite_uniform_conditioning}
The event $C_n$ is finite. Every word in it has length $m=\sum_i n_i$
and the same singleton mass $e^{-Nt}t^m/m!$, where $N=|I|$.
Theorem~\ref{thm:finite_uniform_conditioning} applies to this finite
probability measure whenever the event mass is nonzero.
For $t>0$, that common mass is strictly positive, and $C_n$ is nonempty.
Consequently $\mu_t(C_n)=|C_n|e^{-Nt}t^m/m!>0$.
\end{proof}

\begin{lemma}[Null count events at zero time]
\label{lem:poisson_null_count_conditioning}
\lean{PoissonWord.measure_countEvent_zero_time_of_ne,
PoissonWord.cond_countEvent_zero_time_of_ne}
\leanok
\uses{def:poisson_word_count_vector}
At $t=0$, every $n\ne0$ satisfies
\begin{align}
 \mu_0(C_n)&=0,& \mu_0(\,\cdot\mid C_n)&=0.
 \label{eq:poisson_null_count_conditioning}
\end{align}
The second equality uses the totalized conditional-measure convention in
Lean: conditioning on a null event gives the zero measure.
For $n=0$, the fiber consists only of the empty word and has mass one,
so the conditional law is its point mass.
\end{lemma}
\begin{proof}
\leanok
\uses{lem:poisson_word_normalization, def:poisson_word_count_vector}
At zero time $\mu_0$ is the point mass at the empty word, whose count vector
is zero. Thus $C_n$ has zero mass when $n\ne0$; the definition of conditioning
then gives the zero measure. A word with zero total count has length zero.
\end{proof}

\begin{definition}[Count-and-order sampler]
\label{def:poisson_counts_then_uniform}
\lean{PoissonWord.countsThenUniform}
\leanok
\uses{def:poisson_word_count_vector, lem:poisson_word_count_fiber}
Put $\nu_t=\bigotimes_{i\in I}\operatorname{Poisson}(t)$ and define
\begin{align}
 \widehat\mu_t
   &=\sum_{n\in\N^I}\nu_t(\{n\})\operatorname{Unif}(C_n).
 \label{eq:poisson_counts_then_uniform}
\end{align}
This samples independent counts and then a uniform compatible word.
Every $C_n$ is finite and nonempty, including at $t=0$.
\end{definition}

\begin{theorem}[Reconstruction of the word law]
\label{thm:poisson_counts_then_uniform}
\lean{PoissonWord.countsThenUniform_eq_measure,
PoissonWord.instIsProbabilityMeasureCountsThenUniform}
\leanok
\uses{def:poisson_counts_then_uniform, def:poisson_word_measure}
For every $t\ge0$,
\begin{align}
 \widehat\mu_t&=\mu_t.
 \label{eq:poisson_counts_then_uniform_eq}
\end{align}
In particular, the count-and-order sampler is a probability measure.
\end{theorem}
\begin{proof}
\leanok
\uses{thm:poisson_independent_counts, thm:poisson_conditional_order,
lem:poisson_word_normalization}
Fix a word $w$. In (\ref{eq:poisson_counts_then_uniform}), only
$n=C(w)$ can contribute to its singleton mass: all other uniform measures
are supported on disjoint fibers. The count law gives
$\nu_t(\{n\})=\mu_t(C_n)$.
If this mass is positive, (\ref{eq:poisson_conditional_order}) shows that
its contribution is $\mu_t(C_n)\mu_t(\{w\}\mid C_n)=\mu_t(\{w\})$.
If the fiber has zero mass, both the contribution and $\mu_t(\{w\})$
vanish. Thus the two measures agree on every singleton of the countable
word space, proving equality and normalization.
\end{proof}

When $I=\varnothing$, the only count vector is zero, the only word is
empty, and both measures are its point mass for every $t\ge0$.
The reconstruction also holds at $t=0$: impossible count vectors receive
zero mixture weight even though their combinatorial fibers are nonempty.
These results construct no event times, independent-increment process,
pathwise clock coupling, or locality estimate.
